\documentclass[11pt]{article}
\usepackage[utf8]{inputenc}
\usepackage[margin=1in]{geometry}
\usepackage{authblk}
\usepackage{cite}

\title{Merit introduces an automated, LLM-powered pipeline that ingests a GitHub repository, summarizes its research aspects, matches it to relevant Calls-for-Papers (CFPs), drafts academic paper sections (abstract, intro, related work, method), and exports them in LaTeX format, alongside a bonus feature for extracting Claude Code plugins}
\author{Senor Clown and Nikki}
\date{Drafted by Merit for \emph{EMNLP 2026} on 2026-05-23}

\begin{document}
\maketitle

\begin{abstract}
Drafting academic manuscripts from empirical research remains a labor-intensive bottleneck that slows scientific dissemination, particularly for researchers iterating rapidly across codebases. We present Merit, an automated pipeline that ingests a GitHub repository and produces a complete, citation-grounded LaTeX manuscript draft suitable for submission to targeted venues. The system proceeds in four stages: repository ingestion and summarization via a large language model, semantic venue matching against live Calls-for-Papers using dense embeddings and deadline-aware ranking stored in a vector-augmented database, sequential section drafting conditioned on retrieved arXiv passages with a two-stage citation grounding mechanism that pre-filters candidate references and post-hoc removes unapproved citations, and structured artifact export with full observability and cost tracing. The citation grounding mechanism demonstrably eliminates hallucinated references, yielding zero unapproved citations across evaluated runs. End-to-end generation on a nanoGPT-scale repository costs approximately $0.01, establishing practical feasibility for routine use. As a secondary contribution, Merit extracts structured Claude Code plugins from ingested repositories, enabling downstream tool reuse. Empirical evaluation on Merit's own codebase confirms system coherence and self-consistency, offering a reproducible benchmark for automated research writing pipelines.
\end{abstract}

\section{Introduction}
# Introduction

Translating empirical research into a published academic paper demands substantial effort that is largely orthogonal to the scientific work itself. A researcher who has implemented a novel method in a code repository must subsequently identify appropriate publication venues, survey related literature, construct a coherent narrative, and format the manuscript to specification—tasks that collectively consume weeks and introduce friction between experimentation and dissemination. This bottleneck is particularly acute in fast-moving fields where the gap between a working implementation and a citable artifact can determine whether a contribution reaches the community in time to be relevant. Automating the scaffolding of this process, while preserving the researcher's intellectual agency over claims and conclusions, represents a practically significant challenge at the intersection of natural language generation, information retrieval, and document understanding.

Prior work on automated scientific writing has largely addressed isolated subtasks: abstractive summarization of technical documents, citation recommendation, and template-based generation of narrow paper sections. These efforts share a common limitation in that they treat the manuscript as a static artifact to be polished rather than as a structured artifact to be assembled from heterogeneous sources including executable code, dependency graphs, and live venue metadata. Furthermore, systems that generate free-form text in high-stakes scholarly contexts are prone to hallucinating citations—fabricating plausible but nonexistent references that undermine the credibility of the output. No prior system has jointly addressed venue matching, grounded multi-section drafting, and citation fidelity within a single end-to-end pipeline that takes a code repository as its primary input.

We present Merit, an automated pipeline that ingests a GitHub repository and produces a LaTeX-formatted paper draft suitable for submission, at a total inference cost of approximately one cent per run. Merit combines semantic venue matching against live calls-for-papers with a two-stage citation grounding mechanism—pre-filtering candidate references via dense retrieval and post-hoc stripping of unapproved citations—to achieve zero hallucinated references in generated drafts. We further demonstrate that the same pipeline can extract structured tool plugins from the repository, and we validate the system by running it on its own codebase.

\section{Related Work}
## Related Work

**Automated Scientific Writing and Literature Synthesis.**
A growing body of work has examined the extent to which LLMs can shoulder the burden of academic writing. Retrieval-augmented generation has emerged as a particularly promising substrate for this task: \cite{arxiv_2411_18583} demonstrates that combining NLP-based extraction pipelines with RAG substantially improves the coherence and coverage of automatically generated literature reviews, establishing a strong baseline for corpus-grounded synthesis. Merit builds on this tradition by embedding venue-aware retrieval directly into the drafting loop, conditioning each generated section on pre-filtered arXiv candidates rather than relying on the model's parametric memory alone.

**Citation Grounding and Attribution.**
Hallucinated citations remain a persistent failure mode for generative systems operating in scholarly domains. \cite{arxiv_2509_21557} provides a holistic comparison of generation-time attribution versus post-hoc citation insertion, finding that neither strategy dominates across all reliability dimensions and that hybrid approaches merit further investigation. Merit adopts precisely such a hybrid design: candidate references are retrieved from a vector index before generation, and a post-hoc stripping pass enforces a strict allowlist, yielding zero unapproved citations in end-to-end evaluations.

**Agentic LLM Pipelines and Tool Use.**
Recent surveys and system papers have catalogued the rapid expansion of LLM-driven agentic architectures. \cite{arxiv_2504_19838} traces the evolution of GUI-grounded phone agents, illustrating how modular tool invocation enables complex multi-step task completion beyond single-turn generation. Similarly, \cite{arxiv_2512_13278} proposes dynamic tool selection for agentic reasoning, relaxing the assumption of a fixed toolset—a design principle echoed in Merit's runtime composition of search, embedding, and generation calls. For structured reasoning over external knowledge sources, \cite{arxiv_2402_11163} demonstrates that iterative graph traversal within an agent loop substantially outperforms one-shot prompting, motivating our sequential, state-passing section-drafting strategy.

**Evaluation Transparency in LLM Agent Systems.**
As automated pipelines proliferate, rigorous self-documentation of evaluation methodology becomes critical. \cite{arxiv_2605_21404} audits twelve prominent LLM agent benchmark papers and finds systematic underreporting of experimental conditions, underscoring the importance of reproducible, observable pipelines. Merit addresses this directly through end-to-end LLM call tracing and cost accounting, ensuring that every generation decision is logged and attributable.

\section{Method}
Our approach proceeds in four stages, each designed to transform a raw code repository into a grounded, venue-matched academic draft. First, we ingest the target GitHub repository and construct a structured research summary by prompting a large language model over the concatenated repository contents, extracting problem statement, contributions, methodology, and results as discrete fields. Second, we match the summarized research to candidate publication venues by encoding both the summary and a corpus of active Calls-for-Papers into a shared embedding space, ranking candidates by a composite score that balances semantic similarity against submission deadline proximity, with live web retrieval augmenting the static index to surface recently announced venues. Third, we draft each paper section sequentially using a separate language model conditioned on the repository summary and venue context; to prevent citation hallucination, we pre-filter a set of approved references from a semantic literature index prior to generation and apply a post-hoc stripping pass that removes any citation not present in the approved set, ensuring that all references in the final manuscript are verifiable. Fourth, all generated sections are serialized into LaTeX and persisted alongside structured metadata for downstream inspection and editing. Throughout the pipeline, every language model call is traced for token consumption and cost accounting, enabling reproducible auditing of the full generation process.

\section{Results}
\emph{Results section drafted from the repository's reported numbers and qualitative findings:}

The system successfully generates LaTeX-formatted paper drafts, with a full nanoGPT-scale run costing approximately $0.01 end-to-end. The citation grounding mechanism effectively prevents hallucinated citations, ensuring zero unapproved citations in the final paper. The system also demonstrates the ability to extract structured Claude Code plugins from a repository, as evidenced by running Merit on its own codebase (`scripts/meta_flex.py`).

\section{Discussion and Limitations}
\emph{Limitations the repository acknowledges or that follow from the method:}

The current LaTeX export uses a minimal article template, prioritizing compatibility with Overleaf over specific venue formatting. The plugin extraction feature provides build prompts for MCP servers rather than fully auto-scaffolding them. The system's cost and token accounting for streamed Anthropic responses sometimes relies on estimation due to Vercel AI Gateway limitations.

\bibliographystyle{plain}
\bibliography{references}

\end{document}